% Formato de la serie Informes sobre la Productividad en Colombia del CJC.
\usepackage[margin=1in]{geometry}
\usepackage[utf8]{inputenc}
\usepackage[T1]{fontenc}
\usepackage[spanish]{babel}
\usepackage{graphicx}
\usepackage{xcolor}
\usepackage{mdframed}
\usepackage{setspace}
\usepackage{titlesec}
\usepackage{fancyhdr}
\usepackage{hyperref}
\usepackage{enumitem}
\setstretch{1.15}
\setlength{\parindent}{0pt}
\setlength{\parskip}{8pt}
\titleformat{\section}{\large\bfseries}{\thesection.}{0.5em}{}
  [\vspace{2pt}\hrule\vspace{4pt}]
\titlespacing{\section}{0pt}{18pt}{8pt}
\titleformat{\subsection}{\normalsize\bfseries}{}{0pt}{}
\titlespacing{\subsection}{0pt}{14pt}{6pt}
\hypersetup{colorlinks=true,linkcolor=blue!60!black,citecolor=blue!60!black,urlcolor=blue!60!black,
  pdftitle={Productividad innovación y financiación empresarial},pdfauthor={Oliver Pardo}}
\graphicspath{{Paper/figures/}}
\newcommand{\reportseries}{Informes sobre la Productividad en Colombia}
\newcommand{\reportseriesdesc}{Evidencia para entender y fomentar la productividad}
% Etiqueta de edición. Sustituir cuando se asigne un número a la publicación.
\newcommand{\reportnumber}{Propuestas para el PND}
\newcommand{\reporttitle}{Productividad, innovación y financiación empresarial}
\newcommand{\reportsubtitle}{Diez propuestas para el Plan Nacional de Desarrollo}
\newcommand{\reportauthor}{Oliver Pardo}
\newcommand{\reportaffiliation}{Centro Javeriano de Competitividad\\Pontificia Universidad Javeriana}
\newcommand{\reportdate}{26 de octubre de 2026}
\pagestyle{fancy}
\fancyhf{}
\fancyhead[L]{\footnotesize Centro Javeriano de Competitividad}
\fancyhead[R]{\footnotesize \reportnumber}
\cfoot{\thepage}
\setlength{\headheight}{14pt}
\renewcommand{\headrulewidth}{0.4pt}
\renewcommand{\refname}{Referencias}
